% Scalability analysis (single USL curve), limiting pipeline stages with mitigations, and eBPF probes.
% Generated inputs: usl_text.tex (usl/make_usl_tex.py), usl_single_fit.png (usl/usl_single_fit.py),
% ebpf_summary.tex (ebpf/summarize.py). Everything marked TODO needs the measurements.

\section{Scalability Analysis}
\label{sec:usl}

\subsection{One USL curve}

Following the feedback on Stage I we fit a single curve instead of one per ensemble size, and we make it a well-supported one:
one configuration, many load levels, replicated runs in random order, on dedicated machines.
The effect of the ensemble size is not read from separate fits but measured as factor A of the factorial
experiment (Section~\ref{sec:design}).

\input{usl_text.tex}

\begin{figure}[h]
  \centering
  \IfFileExists{usl_single_fit.png}{\includegraphics[width=0.8\textwidth]{usl_single_fit.png}}{\textbf{[TODO: usl\_single\_fit.png]}}
  \caption{Measured throughput (dots: runs; black: mean $\pm$ one standard deviation), the USL fit with its 95\% bootstrap band, and the
    contention-only (Amdahl) model. The shaded region on the right was held out when fitting for the validation.}
  \label{fig:usl}
\end{figure}

\subsection{Which parts of the pipeline limit scalability}

Each term of the model corresponds to a part of the write pipeline of Section~2 of Stage I
(Table~\ref{tab:limits}). Reads are different: a read is answered by the server the client is connected to, with no
coordination, so reads are limited by the CPU of that server and not by the leader.

\begin{table}[h]
  \centering
  \caption{USL terms, the pipeline stage each one points to, the evidence of this work and a mitigation.}
  \label{tab:limits}
  \footnotesize
  \begin{tabular}{@{}p{1.0cm}p{3.3cm}p{3.4cm}p{3.8cm}@{}}
    \toprule
    Term & Pipeline stage & Evidence & Mitigation \\
    \midrule
    $\alpha$ & Serialized request pipeline of the leader: request preparation, then the log-sync stage
      (append and \texttt{fsync} of the transaction log) and the commit stage
      & \textbf{[TODO: effect of D (forceSync) and F (log disk) in the factorial; fsync probe: calls and p99]}
      & Faster, dedicated log device; batching of log writes (group commit); fewer or larger writes by the
        clients (multi-operations) \\
    $\beta$ & Quorum acknowledgment: the leader must exchange messages with every follower and wait for a
      majority; cost grows with the ensemble
      & \textbf{[TODO: effect of A and A$\times$B; TCP retransmit probe; coherency term of the fit]}
      & Keep the voting ensemble small; add \emph{observers} (non-voting servers) to scale reads without enlarging
        the quorum; hierarchical quorums or weights across zones \\
    $\gamma$ & Cost of one request in isolation: network round trips, serialization, payload bytes, the fsync
      latency for a single write
      & \textbf{[TODO: effect of E (payload) and of the one-client latencies]}
      & Asynchronous clients that pipeline several requests per session; smaller payloads; co-located clients \\
    \bottomrule
  \end{tabular}
\end{table}

\textbf{[TODO: one paragraph that names the stage that dominates in the measurements, with the numbers, and
states which mitigation the data support. If forceSync (D) has no effect, say that the log-sync stage is not the
bottleneck and look at the CPU of the leader instead (run-queue probe).]}

\section{eBPF Probes}
\label{sec:ebpf}

We wrote four probes that observe the server's operating system while the benchmark runs, each aimed at one of the
suspected limits of Table~\ref{tab:limits}.
\begin{itemize}
  \item \textbf{fsync latency} (custom bpftrace program on the tracepoints \texttt{sys\_enter/exit\_fsync} and
    \texttt{fdatasync}, filtered to the ZooKeeper process): the time spent in the log-sync stage and the number of
    calls. With \texttt{forceSync=no} it must show no calls, which also verifies that factor D was applied.
  \item \textbf{Block I/O latency} (bcc \texttt{biolatency}): whether the log disk itself is slow underneath the
    fsync; compared between pd-balanced and pd-ssd.
  \item \textbf{Run-queue latency} (bcc \texttt{runqlat} on the ZooKeeper process): whether the server is waiting
    for a CPU, which would add to $\alpha$ independently of the disk.
  \item \textbf{TCP retransmissions} (custom bpftrace program on \texttt{tcp:tcp\_retransmit\_skb}): whether packet loss
    between the servers adds retransmission timeouts to the quorum wait, which would show up in $\beta$.
\end{itemize}
Probes have a small overhead, so they were run in separate runs of the saturated configuration, at both levels of
\texttt{forceSync} and of the log disk, and never during the runs of the factorial experiment.

\begin{table}[h]
  \centering
  \caption{eBPF measurements on the leader during a 40\,s saturated run; percentiles are upper bounds of
    histogram buckets.}
  \label{tab:ebpf}
  \footnotesize
  \IfFileExists{ebpf_summary.tex}{\input{ebpf_summary.tex}}{\textbf{[TODO: run ebpf/summarize.py --tex report/ebpf\_summary.tex]}}
\end{table}

\textbf{[TODO: interpretation. What fraction of the write latency is the fsync (compare its p50 with the write latency
of the benchmark)? Does the fsync latency change with the disk type as the block-I/O histogram predicts? Is the
run queue empty? Are there retransmissions? State which hypotheses of Section~\ref{sec:design} these results
support.]}
